\clearpage
\fancyhead[L]{\small CONNAISSANCE DES TEMS.}
\newcommand{\ObservedVenusFraction}{\hbox{\vrule width0pt height7.600000bp depth0pt\special{pdf:bcontent}\special{pdf:code q 0 g 0.205405 4.313514 m 0.000000 4.108108 l 0.000000 3.902703 l 0.205405 3.697297 l 0.821622 3.697297 l 1.027027 3.902703 l 1.232432 3.902703 l 1.437838 3.697297 l 1.437838 3.286486 l 1.232432 3.081081 l 1.027027 3.081081 l 0.821622 2.875676 l 0.821622 2.054054 l 1.232432 1.643243 l 1.232432 1.232432 l 1.027027 1.027027 l 1.027027 0.205405 l 1.232432 0.000000 l 1.437838 0.000000 l 1.643243 0.205405 l 1.848649 0.205405 l 2.054054 0.000000 l 2.464865 0.000000 l 2.670270 0.205405 l 2.875676 0.205405 l 3.697297 1.027027 l 3.697297 2.670270 l 3.081081 3.286486 l 3.081081 3.697297 l 3.286486 3.902703 l 3.697297 3.902703 l 3.902703 3.697297 l 4.724324 3.697297 l 4.929730 3.902703 l 4.929730 4.108108 l 4.724324 4.313514 l 3.286486 4.313514 l 3.081081 4.108108 l 1.027027 4.108108 l 0.821622 4.313514 l h 2.259459 7.600000 m 1.027027 6.367568 l 1.027027 6.162162 l 1.232432 5.956757 l 1.437838 5.956757 l 1.848649 5.545946 l 1.848649 5.135135 l 2.259459 4.724324 l 2.875676 4.724324 l 3.081081 4.929730 l 3.081081 5.135135 l 2.875676 5.340541 l 2.875676 7.394595 l 2.670270 7.600000 l h f* Q}\special{pdf:econtent}\kern4.929730bp}}

\hypertarget{native-586-lower}{}
\begin{center}\small Göttingische gelehrte Anzeigen. Stück 40. Seiten 385..396. 11. März 1815.\\[3pt]\rule{.72\linewidth}{.3pt}\end{center}

\emph{Connaissance des tems, ou des mouvemens célestes, à l'usage des astronomes et des navigateurs pour
l'an} 1816; \emph{ publiée par le bureau des longitudes.} Paris. Bei der Witwe Courcier. November 1813. 360 Seiten im Großoktavformat.

Wir übergehen den astronomischen Kalender und die regelmäßig erscheinenden Artikel mit Stillschweigen und besprechen nur die \emph{Zusätze}, die seit mehreren Jahren fast ganz gefehlt hatten und im vorliegenden Jahrgang wieder zahlreich und interessant sind. Zuerst kommt ein bemerkenswerter Aufsatz über die Kometen von Graf \textsc{Laplace}. Dieser große Geometer spricht sich hier für die Hypothese aus, nach der jene wunderbaren Himmelskörper durch Verdichtung desselben Stoffs entstehen, der die Nebelsterne bildet und im ganzen Weltraum so allgemein verbreitet zu sein scheint. Demnach stünden die Kometen zu unserem Sonnensystem nur in derselben zufälligen Beziehung wie die \SourceLeaf{587}Meteoriten nach Ansicht einiger Naturforscher zur Erde. Gelangen sie so weit in den Bereich unserer Sonne, dass deren Anziehung überwiegt, so werden sie gezwungen, elliptische oder hyperbolische Bahnen um sie zu beschreiben. Daraus erklärt sich von selbst die große Vielfalt oder vielmehr die völlige Regellosigkeit in der \emph{Lage} der beobachteten Kometenbahnen. Nur scheint sich ein Einwand gegen diese Hypothese daraus zu ergeben, dass wir unter allen bisher berechneten Bahnen, deren Zahl schon weit über hundert liegt, noch keine einzige eindeutig hyperbolische erkannt haben. Denn die Art des Kegelschnitts hängt von der ursprünglichen Geschwindigkeit des Kometen an der Stelle ab, an der er in die Sphäre der überwiegenden Sonnenanziehung eintritt, und es scheint sich ebenso gut eine Hyperbel wie eine Ellipse ergeben zu können. Der Aufsatz soll nun mit den Mitteln der strengen Wahrscheinlichkeitsrechnung beweisen, dass unter den Bahnen aller Kometen, die uns nahe genug kommen, um beobachtet werden zu können, nur ein äußerst geringer Teil Hyperbeln mit einer so kleinen Achse sind, dass sie anhand der Beobachtungen als Hyperbeln erkannt werden können. Für das Verhältnis der Wahrscheinlichkeit einer solchen Hyperbel, deren große Halbachse 100 Radien der Erdbahn nicht übersteigt, zur Wahrscheinlichkeit der übrigen Fälle, nämlich einer Hyperbel mit größerer Achse, einer Parabel oder einer Ellipse, findet \textsc{Laplace} den Ausdruck (I)
\[1 : \frac{\pi-2}{10}\BareRoot{\frac{r(r+200)}{2D}} - 1\]
Dabei bezeichnet $\pi$ den halben Kreisumfang bei einem Radius von 1, $D$ den größten Perihelabstand für noch beobachtbare Kometen und $r$ den Radius der Sphäre der überwiegenden Sonnenanziehung. Setzt er hier $r = 100000$, $D = 2$, Voraussetzungen, die man ohne Weiteres gelten lassen kann, so ergibt sich, dass man 56 gegen 1 darauf wetten könnte, dass unter hundert beobachteten Kometenbahnen keine erkennbare Hyperbel vorkommt. Dabei wird vorausgesetzt, dass alle Perihelabstände von 0 bis 2 gleich möglich sind. Nimmt man an, dass die Wahrscheinlichkeit mit zunehmendem Perihelabstand immer weiter abnimmt, so wird das obige Verhältnis noch größer. Unter einer sehr plausiblen Hypothese für die abnehmende Wahrscheinlichkeit bei zunehmendem Perihelabstand berechnet \textsc{Laplace}, dass man 82 gegen eins darauf wetten könnte, dass unter hundert beobachteten Bahnen keine erkennbare Hyperbel vorkommt.

Die Bedeutung des Gegenstands und der Name des Verfassers rechtfertigen es, bei einem Aufsatz von nur wenigen Seiten länger zu verweilen. Das Vorstehende sollte nur historisch über seinen Inhalt und Zweck berichten. Was wir noch hinzufügen, ist nur für diejenigen unserer mathematischen Leser bestimmt, die den Aufsatz selbst studieren und vor sich haben. So gern wir der Hypothese selbst den Vorzug vor allen anderen geben, die über diese rätselhaften Himmelskörper aufgestellt worden sind, namentlich vor der von \textsc{Lagrange} in der Connaissance des tems für 1814, müssen wir doch zum mathematischen Teil des Aufsatzes zwei Anmerkungen machen. Erstens scheint uns \textsc{Laplace} auf S. 216 voreilig zu urteilen, wenn er sagt `la plus petite valeur de $V$ est celle qui rend nulle la quantité renfermée sous le radical précédent.' Der kleinste Wert von $V$ ist vielmehr 0. Der Ausdruck (II)
\[1{}^{\displaystyle -}\frac{\BareRoot{(1-\frac{D}{r})}}{rV}\cdot\BareRoot{(rrVV(1+\frac{D}{r})-2D)}\]
für die Wahrscheinlichkeit, dass der Perihelabstand zwischen 0 und $2D$ liegt, gilt noch nicht ab $V=0$. Er gilt erst ab dem Wert von $V$, den \textsc{Laplace} als kleinsten bezeichnet. Für kleinere Werte von $V$ ist diese Wahrscheinlichkeit nicht imaginär, wie es Ausdruck (II) ergeben würde, sondern \SourceLeaf{588}$= 1$, das heißt \emph{Gewissheit}. Diese kleineren Werte von $V$ hat \textsc{Laplace} nach unserer Auffassung zu Unrecht außer Acht gelassen. Zu dem von \textsc{Laplace} entwickelten Integral muss also noch $1\times\int dV$, genommen von $V = 0$ bis $V = \BareRoot{\dfrac{2D}{1+\dfrac{D}{r}}}$, das heißt eben die Größe $\BareRoot{\dfrac{2D}{1+\dfrac{D}{r}}}$ selbst, hinzugefügt werden. Behielten wir \textsc{Laplace}s Entwicklungen im Übrigen unverändert bei, so erhielten wir statt des oben angegebenen Ausdrucks (I) das folgende Verhältnis.
\[1 : \frac{\pi}{10}\BareRoot{\frac{r(r+200)}{2D}} - 1\]
Damit könnte man statt 56 tatsächlich 157 gegen 1 wetten, was für die Hypothese noch viel günstiger wäre. Wir haben aber noch eine zweite Bemerkung zu machen, die zu einem ganz anderen Ergebnis führt. \textsc{Laplace} entwickelt auf S.\,217 den von ihm auf S.\,216 gefundenen Integralausdruck (b) in eine Reihe und erhält (III)
\[\frac{(\pi-2)\BareRoot{2D}}{2r} - \frac{D}{ir\BareRoot{r}}\]
Dafür wäre also nach unserer ersten Bemerkung
\[
\frac{\pi\BareRoot{2D}}{2r} - \frac{D}{ir\BareRoot{r}}
\]
einzusetzen. Diese Ausdrücke nähern sich also endlichen Grenzen, wenn $i$ ins Unendliche wächst. Das kann nicht sein, wie man sich anhand des exakten Integrals oder sogar schon anhand des Differentials vor der Integration leicht überzeugen kann. Führt man die Entwicklung weiter, so findet man tatsächlich, dass (IV) noch hinzuzufügen ist.
\[+\,\frac{iDD}{2rr\BareRoot{r}}\]
Weil \textsc{Laplace} diesen Teil außer Acht ließ, konnte er sein Endergebnis von $U$ unabhängig machen, indem er diese Größe als unendlich betrachtete (was \textsc{Laplace} eigentlich stillschweigend voraussetzt, indem er auf S.\,217 $i$ unendlich groß nimmt). Wollte man aber dieselbe Voraussetzung, dass für die ursprünglichen Geschwindigkeiten jeder Wert ohne Einschränkung gleich möglich sei --- weiterhin gelten lassen und dabei das Glied (IV) berücksichtigen, so würde das Endergebnis nun ganz anders ausfallen. Es wäre unendlich wenig wahrscheinlich, dass sich unter einer endlichen Anzahl beobachteter Kometenbahnen irgendeine Ellipse, Parabel oder einer Parabel nahekommende Hyperbel befände. Vielmehr wären sie aller Wahrscheinlichkeit nach sämtlich von geraden Linien nicht zu unterscheidende Hyperbeln. Offenbar ist aber die Voraussetzung selbst unzulässig, weil sie jeder endlichen Geschwindigkeit unendlich wenig Wahrscheinlichkeit ließe. Man \emph{muss} für $U$ eine endliche Grenze annehmen. Da das Endergebnis dann ganz von $U$ abhängig bleibt und wir keinen Anhaltspunkt haben, darüber etwas festzulegen, bleibt die Aufgabe eigentlich unlösbar. Kann die Wahrscheinlichkeitsrechnung auch keinen entscheidenden Beweis \emph{für} die Hypothese liefern, so entscheidet sie doch gerade wegen unserer Unkenntnis der Grenze $U$ auch durchaus nichts gegen die Hypothese. Die Rechnung selbst zeigt, dass erkennbar hyperbolische Bahnen immer sehr selten bleiben müssen, solange man keine \emph{sehr} großen Werte für $U$ annimmt. Eine Überschlagsrechnung ergibt, dass wir mindestens $U = 1083$ annehmen müssen (die Einheit ist die mittlere Geschwindigkeit der Erde auf ihrer Bahn), um erwarten zu dürfen, unter hundert Kometenbahnen eine zu finden, die wenigstens eindeutig hyperbolisch ist.

\SourceLeaf{589}Bei der Besprechung der übrigen Aufsätze dürfen wir uns kürzer fassen. Der gelehrte Herausgeber der C. d. T. untersucht im nächsten Aufsatz die Frage, ob \textsc{Hipparch} in Alexandria beobachtet hat, und macht es wahrscheinlich, dass dieser größte Astronom des Altertums entweder nie in Alexandria gewesen ist oder sich dort höchstens vorübergehend aufgehalten hat. --- Vom selben Verfasser über die Azimutbeobachtungen, die \textsc{Tycho} in Uranienburg angestellt hat (sie sind freilich ziemlich fehlerhaft, stammen aber nur aus der ersten Zeit seines Aufenthalts auf der Insel Hveen und berechtigen uns nicht, die richtige Aufstellung seiner Meridianinstrumente anzuzweifeln). --- Bemerkenswerte Auszüge aus der Handschrift von \textsc{Ptolemäus}' Optik, die sich in der Pariser Bibliothek befindet. --- Auszug aus \textsc{Don Joseph Rodriguez}' Abhandlung über die englische Gradmessung mit einigen Anmerkungen. Herr \textsc{Rodriguez} wollte zeigen, dass die teilweisen Abweichungen bei den verschiedenen Gradmessungen uns noch nicht dazu berechtigen, eine unregelmäßige Gestalt des Erdkörpers als erwiesen anzusehen. Herr \textsc{Delambre} bemerkt dagegen zu Recht, dass auch das Gegenteil nicht bewiesen ist. Doch was dieser Astronom zu \textsc{Mechain}s Beobachtungen in Barcellona und Montjouy hinzufügt, scheint uns deren Ergebnis keineswegs zu entkräften. Die mit den Kreisen von \textsc{Lenoir} bestimmten Polhöhen weisen demnach noch Unsicherheiten von mehreren Sekunden auf. Um diese bemerkenswerten Beobachtungen, die Anlass zu manchem Hin und Her gegeben haben, angemessen würdigen zu können, muss man die Ergebnisse eigentlich folgendermaßen vergleichen. Der Meridianbogen zwischen Montjouy und Barcellona, für den die geodätischen Messungen $=59''53$ lieferten, ergibt sich nach den astronomischen Beobachtungen
\begin{center}
\fontsize{9}{11}\selectfont\setlength{\tabcolsep}{3pt}\begin{tabular}{@{}l c r@{}}
aus $\beta$ der Zwillinge & $.\;.\;.\;.\;.\;.\;.\;.\;.\;.\;.\;.\;.\;.\;.\;.$ & $57''20$\\
$\zeta$ Großer Bär obere Kulmination & $.\;.\;.\;.\;.\;.\;.\;.$ & $61.59$\\
$\zeta$ Großer Bär untere Kulmination & $.\;.\;.\;.\;.\;.\;.\;.$ & $64.66$\\
$\beta$ Kleiner Bär obere Kulmination & $.\;.\;.\;.\;.\;.\;.\;.$ & $62.75$\\
$\beta$ Kleiner Bär untere Kulmination & $.\;.\;.\;.\;.\;.\;.\;.$ & $63.39$\\
$\alpha$ Kleiner Bär obere Kulmination & $.\;.\;.\;.\;.\;.\;.\;.$ & $62.42$\\
$\alpha$ Kleiner Bär untere Kulmination & $.\;.\;.\;.\;.\;.\;.\;.$ & $63.12$\\
\end{tabular}
\end{center}
So finden wir hier unabhängig von allen örtlichen Anziehungen Unterschiede von bis zu $7''46$ zwischen den Ergebnissen für verschiedene Sterne, die sich nur der Beobachtung und dem Instrument zuschreiben lassen. --- Noch einiges zu den bekannten Parallaxenformeln von Herrn \textsc{Olbers}, von Herrn \textsc{Delambre}, und vom selben Verfasser eine Besprechung seines Abrégé de l'astronomie und der Übersetzung des Almagests von Halma. --- Verbesserte Elemente der Bewegungen des Jupiter von Herrn \textsc{Bouvard}, nach denen wir bald neue Tafeln erwarten dürfen. --- Verschiedenes zur Höhenmessung mit dem Barometer von Herrn \textsc{Prony}. Eine hier beschriebene neue Methode, die Genauigkeit der Ablesung zu erhöhen, verdient Nachahmung. In mehreren hier angeführten Beobachtungsreihen weichen die äußersten Ergebnisse für den Höhenunterschied im Mittel nur um 2 bis 3 Meter voneinander ab. --- Über die Pendelversuche, die spanische Seefahrer an vielen Punkten der Erde angestellt haben, von Herrn \textsc{Mathieu}. --- Über die neuen Mondtafeln von Herrn \textsc{Burckhardt}. Was wir uns bei der Besprechung dieser vortrefflichen Tafeln gewünscht hatten, nämlich die vollständigen Einzelergebnisse des Vergleichs mit den Beobachtungen, finden wir hier. Bemerkenswert scheint uns, dass die negativen Breitenfehler merklich überwiegen. Wir verstehen das so (ausdrücklich steht es dort allerdings nicht), dass die Tafeln die Breite im Durchschnitt um $2''29$ weiter südlich angeben als die Beobachtungen. Das deutet darauf hin, dass die Rechnungselemente geringfügig korrigiert werden müssen. Wie viel davon aber etwa auf die Parallaxenkonstante, den \SourceLeaf{590}Kollimationsfehler, die Refraktion oder die Polhöhe entfällt, wird schwer zu entscheiden sein. --- Über die Massen der Planeten, vom selben Verfasser. Zur Vorbereitung neuer Tafeln für Mars, Venus und Merkur, die wir, wie wir hieraus sehen, von diesem vortrefflichen Astronomen erwarten dürfen, hat er die Elemente der Erdbewegung anhand von beinahe 4000 Beobachtungen erneut untersucht. Hier erhalten wir die wichtigsten Ergebnisse dieser Arbeit. Bemerkenswert ist, dass sich daraus eine Verringerung der Venusmasse um ungefähr $\ObservedVenusFraction$ ergibt, während Herr \textsc{von Lindenau} aus den Bewegungen des Merkur eine ungefähr ebenso große Zunahme gefunden hat. --- Bemerkungen über die Vervielfältigungskreise von Herrn \textsc{Arago}. Alle Astronomen kennen die Bedenken, die Herr \textsc{von Zach} gegen die Kreise mit fester Libelle an der Achse erhoben hat. Herr \textsc{Arago} sucht sie in diesem Aufsatz auszuräumen. Theoretisch erscheine es zwar begründet, dass die Handgriffe bei den geraden Beobachtungen die Verbindung des Kreises mit der Säule geringfügig verändern könnten, ohne dass die Libelle dies anzeigt. Man müsse darüber aber die Instrumentenmacher selbst als Schiedsrichter befragen. Sie glauben nicht, dass es zu einer merklichen Verschiebung kommen könne, sofern nur die Klammer die richtige Größe habe und sich die Alhidade auf dem Limbus leichtgängig bewege. (Die Bedenken betreffen weniger die Klammer als vielmehr die Stellschraube, die letztlich allein den Kreis mit der Säule verbindet. Uns scheint, dass sich \emph{dieser} Zweifel klären ließe, wenn man bei einigen Beobachtungsreihen ein und desselben Sterns das Fernrohr bei den geraden Beobachtungen stets wie gewöhnlich durch den Zenit auf den Stern zurückführte, bei anderen Reihen dagegen mit einer Bewegung in entgegengesetzter Richtung durch den Nadir. Bewirkte die Bewegung des Fernrohrs tatsächlich eine merkliche Verschiebung des Kreises selbst, so müsste man erwarten, dass die Ergebnisse der ersten Reihen von denen der anderen stets in derselben Richtung abweichen.) Die Erfahrungen, die Herr \textsc{von Zach} gegen Kreise mit fester Libelle anführt, seien hier deshalb nicht entscheidend, weil sie von anderen Beobachtern stammten, die so vortrefflich übereinstimmenden Beobachtungen an Kreisen mit beweglicher Libelle hingegen von Herrn \textsc{von Zach} selbst. Herr \textsc{Arago} verteidigt ferner die von Herrn \textsc{von Zach} missbilligte Beobachtungsmethode. Bei den geraden Beobachtungen lässt man den Stand der Säule selbst unverändert, zeichnet stets nur die Stellung der Blase zu den Teilstrichen der Skala auf und korrigiert entsprechend das Endergebnis. Herr \textsc{Arago} hat den Wert der Teilstriche auf der Libellenskala des dreifüßigen Kreises von \textsc{Reichenbach}, den die Pariser Sternwarte als Geschenk des Grafen \textsc{Laplace} besitzt, zweimal bei sehr unterschiedlichen Temperaturen untersucht und fast genau übereinstimmende Ergebnisse erhalten. Damit ist allerdings \emph{der} Zweifel noch nicht ausgeräumt, ob man berechtigt sei. das Innere der Libellenröhre für so durchweg gleichförmig gekrümmt zu halten, dass \emph{jedem einzelnen} Teilstrich der Skala genau gleiche Werte zukommen. Bedenkt man, dass hier schon ganz außerordentlich kleine Abweichungen von der regelmäßigen Krümmung genügen, um Unterschiede von mehreren Sekunden hervorzurufen, so wird man den Zweifel, ob auch der geschickteste Instrumentenmacher solche Abweichungen ganz vermeiden kann, nicht für unbegründet halten. Solange darüber keine völlige Gewissheit besteht, würden wir es jedenfalls für sicherer halten, bei jeder geraden Beobachtung durch Verstellen der Säule die Luftblase genau an die Stelle zurückzubringen, an der sie bei der vorhergehenden ungeraden Beobachtung stand.

Herr \textsc{Arago} kommt auch noch auf die bemerkenswerte Behauptung von Herrn \textsc{von Zach} zu sprechen, dass bei allen Beobachtungen mit Repetitionskreisen hinsichtlich der absoluten Genauigkeit der Ergebnisse immer noch eine Unsicherheit von mehreren Sekunden bleibt, so gut sie auch untereinander übereinstimmen mögen. Denn \emph{verschiedene} Kreise liefern unterschiedliche Ergebnisse. Da dieser berühmte Beobachter selbst die Ursache dieser Unterschiede noch als Rätsel bezeichnet hat, könnte \SourceLeaf{591}ein anderer Astronom, der nicht einmal alle Begleitumstände kennt, darüber noch weniger entscheiden. Herr \textsc{Arago} führt hier jedoch einige ähnliche eigene Erfahrungen an und glaubt deren Ursache in einem Umstand zu finden, zu dem wir noch einiges hinzufügen müssen. Aus 422 Beobachtungen des Polarsterns mit einem Kreis von \textsc{Fortin} mit fester Libelle, die vier vortrefflich übereinstimmende Reihen bildeten, ergab sich für die Breite von Formentera $38^\circ\ 39'\ 56''05$. Dann wurde das Objektiv etwas verschoben, um eine deutlichere Sicht zu erhalten. Sieben weitere, wiederum vortrefflich übereinstimmende Reihen ergaben aus 570 Beobachtungen $38^\circ\ 39'\ 53''91$. Ähnliche Unterschiede fanden sich später mit demselben Instrument in Paris, wo die Polhöhe je nach Stellung des Objektivs zwischen $48^\circ\ 50'\ 11''$ und $48^\circ\ 50'\ 15''$ schwankte, obwohl sich bei einer bestimmten Stellung stets äußerst genau übereinstimmende Ergebnisse ergaben. Das Ergebnis hing immer eng mit der Stellung des Objektivs zusammen, wenn auch bei verschiedenen Beobachtern auf unterschiedliche Weise. Herr \textsc{Arago} schließt daraus, dass diese Abweichungen allein von der Art herrühren, wie man den Mittelpunkt des Sterns schätzt. Die Fernrohre zeigen das Bild des Sterns immer viel zu groß, und Herr \textsc{Arago} bezweifelt, dass der Mittelpunkt des scheinbaren Bildes zugleich der Mittelpunkt des wahren Bildes ist. Jeder Beobachter habe seine eigene Art, den Mittelpunkt zu schätzen, und so erkläre sich das Phänomen ganz natürlich. Wir müssen gestehen, dass wir diese Erklärung nicht für stichhaltig halten können. Selbst wenn wir einräumen wollten, dass die beiden Mittelpunkte tatsächlich wie beschrieben auseinanderliegen, müsste doch offenbar bei einer halben Drehung des Fernrohrs um seine Achse der zuvor obere Mittelpunkt zum unteren werden. Genau das muss bei der Beobachtung mit Repetitionskreisen geschehen, bei der das Fernrohr in der geraden und der ungeraden Beobachtung entgegengesetzte Lagen einnimmt, sodass sich die Zielfehler aufheben müssen. Und selbst wenn das Fernrohr ein ganz unregelmäßig geformtes Bild erzeugte, wäre es völlig gleichgültig, welchen Punkt darin man anvisierte, solange man bei ungeraden und geraden Beobachtungen stets ein und denselben Punkt wählte. Darf man die Tatsache, wie sie berichtet wird, als ganz gesichert ansehen, so möchte man fast vermuten, dass die Fäden durch das Verschieben des Objektivs außerhalb seines Brennpunkts gerieten, also eine Parallaxe hatten. Dann kam es freilich darauf an, wie der Beobachter sein Auge gewöhnlich vor das Okular hielt. Überhaupt fällt auf, dass das Fernrohr so eingerichtet war, dass man das Objektiv verschieben musste, um deutlich zu sehen. Interessant sind in diesem Aufsatz auch die Ergebnisse, die der oben erwähnte dreifüßige Kreis von \textsc{Reichenbach} für die Polhöhe aus Zenitdistanzen geliefert hat. Man hätte sich nur noch gewünscht, die Rechnungselemente mitgeteilt zu sehen, damit sich diese Beobachtungen zugleich zur Bestimmung sowohl der Deklination als auch der Rektaszension des Polarsterns verwenden ließen. Vielleicht erhalten wir diese Beobachtungen jedoch ausführlicher im nächsten Band der Connoissance des tems, für den auch eine Abbildung des vortrefflichen Instruments versprochen wird.

